
\section{Background on TinyML Performance Analysis}

One the one hand, as the most immediate limiting resource budget on microcontrollers concern memory limitations, typically in the order of kiloBytes, TinyML performance evaluation typically focuses primarily on metrics measuring memory consumption -- while keeping an eye on execution speed -- as described below.
On the other hand, TinyML performance analysis can be tackled at different granularity levels: at the global model level, or at the operator level, for finer granularity, as described in the following.

\subsection{Performance Metrics}
The considered metrics offer insights into the feasibility, efficiency and resource utilization of offloading model inference burden to low-power devices. By analyzing these metrics, users can make initial decisions regarding model  selection, optimization techniques, and hardware configurations to maximize performance and minimize the resource footprint on low-power devices.

{\bf Memory (RAM) Consumption --}
This metric measures the amount of dynamic memory space (primary RAM) consumed by the model during inference. It reflects the memory footprint of the model activation and is important for low-power devices that have limited memory resources. Efficient memory utilization allows for the deployment of larger and more complex models on such devices.

{\bf Storage (Flash memory) Consumption --}
This metric quantifies the amount of storage space, typically in terms of Flash memory region, required to store the compute instruction and associated parameters. It reflects the model's storage footprint on the low-power device. Minimizing storage consumption allows for accommodating multiple models on the device or orchestrating with other essential applications.

{\bf Computational Latency --}
This metric measures the time consumption  of performing inference for each input sample, either at the model level or at the level of individual operators within the model. It reflects the inference speed of the model on the low-power device and plays crucial role in real-time or latency-sensitive applications. Core clock frequency, cache strategies and communication latency between memory and working core have great impact on this indicator.

{\bf SoC Price --}
This metric considers the cost of the System-on-a-Chip (SoC) used in the low-power device. The price of the SoC affects the overall affordability and feasibility of deploying model in large-scale distributed system. Lower-cost SoCs can make the deployment more accessible and cost-effective.

\subsection{Measurement Granularity}
As for performance analysis of machine learning in other domains, TinyML performance can be measured at different granularity levels:

{\bf Per-Model Evaluation --}
At this coarse level, one measures performance of the model as a whole, i.e. the resource footprint incurred by the execution of the model including all its layers and operators. For example, this allows for evaluating the resource consumption for inference with the production-ready code, on a particular industrial hardware setup.
\iffalse
This approach measures the resource footprint at model level.  It uses the model code generated by AoT executor of uTVM, which is the production version for deployment and allows for simulation of the resource consumption of inference in industrial fields. With user-specific random inputs (uniform distribution by default), user can simulate various data scenarios. By measuring the model-level resource footprint, U-TOE provides insights into the IoT system co-design with other software components. 
\fi

{\bf Per-Operator Evaluation --}
At this level, one measures separately the performance of one or more operators (i.e. one or more components of the model). This per-operator measurement can help identifying specific operators that contribute to performance or inefficiencies, in optimizing the model's efficiency and spotting potential bottlenecks.

\iffalse
U-TOE facilitates the evaluation of computational latency and memory usage for each operator within the model. This evaluation is achieved by leveraging the Remote Procedure Call (RPC) functionality provided by uTVM. U-TOE can then collect performance data for individual operators within the model. Additionally, U-TOE annotates relevant parameters (layer weights and bias etc.) associated with each operator, which enables further profiling and analysis. This per-operator measurement allows users to identify specific operators that contribute to performance or inefficiencies, aiding in optimizing the model's efficiency and identifying potential bottlenecks.
\fi


\iffalse
\subsection{Useful Background on RIOT, uTVM and IoT-Lab}
U-TOE combines the use of RIOT as general-purpose embedded software platform, uTVM as transpiler/compiler of models, and IoT-Lab as open access IoT hardware testbed. Hence, we provide below some additional background on these building blocks.

\emph{RIOT}

\emph{IoT-Lab}

\emph{uTVM}

\fi
